\section{Detailed Mathematical Derivations}
\label{sec:appendix_math}

This section provides complete mathematical derivations and formalisms for the core algorithmic components of the \textbf{Q-SPS} and \textbf{CG-Q-SPS} frameworks, expanding upon the condensed formulations in Section~3.

\subsection{Quantization-Aware Scale Calibration (QASC) Decoupling}
\label{sec:appendix_qasc}

In uniform symmetric quantization, the quantizer maps a real-valued tensor $X$ (either weights $W$ or activations $h$) to a low-bit integer representation $X_q$ over a grid defined by the scale factor $s$. For a $b$-bit signed integer format, the quantization and de-quantization operations are defined as:
\begin{align}
    \text{quant}(X, s, b) &= \text{clip}\left(\left\lfloor \frac{X}{s} \right\rceil, -2^{b-1}, 2^{b-1}-1\right) \\
    \text{dequant}(X_q, s) &= s \times X_q
\end{align}
Let $\tilde{X}(s) = \text{dequant}(\text{quant}(X, s, b), s)$ denote the quantized and subsequently de-quantized tensor. Under post-training quantization, the optimal scale factor $s^*$ is typically found by minimizing the Mean Squared Error (MSE) between the unquantized and quantized tensors:
\begin{equation}
    s^* = \arg\min_{s} \| X - \tilde{X}(s) \|_2^2
\end{equation}

For a low-rank adapter block consisting of a down-projection matrix $A_k \in \mathbb{R}^{r \times D}$ and an up-projection matrix $B_k \in \mathbb{R}^{D \times r}$ with intermediate activation vector $h'_b = A_k h_b \in \mathbb{R}^{r}$, the joint optimization of weight scales $s_A, s_B$ and activation scale $s_h$ to preserve the output of the low-rank block represents a highly non-linear, multi-dimensional joint optimization problem:
\begin{equation}
    \{s_A^*, s_B^*, s_h^*\} = \arg\min_{\{s_A, s_B, s_h\}} \left\| B_k A_k h_b - \tilde{B}_k(s_B) \tilde{A}_k(s_A) \tilde{h}_b(s_h) \right\|_2^2
\end{equation}
The computational complexity of a grid search or gradient-free optimization over this joint 3D space scales as $O(N^3)$, where $N$ is the number of candidate scale steps per dimension, making it highly impractical for edge deployments.

\textbf{QASC Sequentially Decoupled Solution:} We sequentially decouple this optimization into two independent, linear-complexity ($O(N)$) search phases, leveraging the feed-forward structure of the low-rank projection chain:

\paragraph{Phase 1: Down-Projection and Activation Scale Calibration}
We group the activation $h_b$ and the down-projection adapter weights $A_k$ to optimize $s_A$ and $s_h$ sequentially. Since activations exhibit dynamic, sample-dependent range variations (especially under visual feature outliers), we employ dynamic max-abs scaling for intermediate activations during serving. Thus, we first calibrate the weight scale $s_A^*$ of the down-projection adapter offline over a calibration set by minimizing the representation reconstruction error:
\begin{equation}
    s_A^* = \arg\min_{s_A} \left\| A_k h_b - \tilde{A}_k(s_A) h_b \right\|_2^2
\end{equation}
Because the down-projection maps the high-dimensional feature $h_b \in \mathbb{R}^D$ to a very compact bottleneck space ($r=8$), the scale factor $s_A^*$ is optimized by evaluating a discrete set of candidates $s_A \in [\alpha \cdot \max(|A_k|), \max(|A_k|)]$ where $\alpha \in [0.5, 1.0]$. 

\paragraph{Phase 2: Up-Projection Scale Calibration}
Once $s_A^*$ is fixed, the intermediate representations $h'_b$ are computed. We then isolate the up-projection scale factor $s_B^*$ by minimizing the output reconstruction error under the pre-calculated, quantized intermediate representations:
\begin{equation}
    s_B^* = \arg\min_{s_B} \left\| B_k h'_b - \tilde{B}_k(s_B) \tilde{h}'_b \right\|_2^2
\end{equation}
where $\tilde{h}'_b = \text{quant}(h'_b, s_{h'}, 8)$ and $s_{h'}$ is the dynamic, per-sample activation scale factor computed as $s_{h'} = \max(|h'_b|) / 127$. By sequencing the optimization as $s_A^* \rightarrow s_{h'} \rightarrow s_B^*$, the joint $O(N^3)$ search space collapses into two independent $O(N)$ line-searches. This allows our post-training calibration to run in less than 0.1 seconds per adapter on a single CPU core while successfully minimizing the compounding quantization rounding noise and clipping errors.

\subsection{Symmetric Manifold De-Entangling (SMD) via L{\"o}wdin Orthogonalization}
\label{sec:appendix_smd}

When multiple task-specific experts are registered, their routing templates (representation centroids) are calculated over the calibration set. Let $C = [c_1, c_2, \dots, c_K]^T \in \mathbb{R}^{K \times D}$ denote the centroid template matrix, where each row $c_k \in \mathbb{R}^D$ represents the unit-norm centroid vector for task $k$.

Under extreme representation entanglement (characterized by an entanglement factor $\epsilon \approx 0.8$, meaning centroids share significant cosine similarity), the raw centroids $c_k$ are highly non-orthogonal, i.e., $c_i \cdot c_j \approx \epsilon$ for $i \neq j$. Performing nearest-centroid routing over these raw coordinates leads to highly diffused routing weights and high-frequency routing flicker.

While Gram-Schmidt Cross-Centroid Orthogonalization (GS-CCO) orthogonalizes the templates, it is inherently asymmetric and order-dependent. Specifically, if centroids are processed in the order $1 \rightarrow 2 \rightarrow \dots \rightarrow K$, the first centroid $c_1$ remains unchanged, while subsequent centroids are projected onto the orthogonal complement of the previous ones. This causes severe, asymmetric representational distortion and elevates routing flicker on the later-processed task paths.

\textbf{L{\"o}wdin Symmetric Manifold De-Entangling:} To resolve this asymmetry, we employ L{\"o}wdin Symmetric Orthogonalization. SMD treats all task templates symmetrically, finding the orthogonalized template matrix $C_{\text{SMD}} \in \mathbb{R}^{K \times D}$ that is closest to the original centroid matrix $C$ in the least-squares sense:
\begin{equation}
    C_{\text{SMD}} = \arg\min_{U \in \mathbb{R}^{K \times D}} \| U - C \|_F^2 \quad \text{subject to} \quad U U^T = I_{K \times K}
\end{equation}
where $\| \cdot \|_F$ is the Frobenius norm.

The closed-form solution is derived as follows:
\begin{enumerate}
    \item Compute the $K \times K$ symmetric, positive-definite overlap (Gram) matrix $S$:
    \begin{equation}
        S = C C^T \in \mathbb{R}^{K \times K}
    \end{equation}
    where the entry $S_{i, j} = c_i \cdot c_j$ represents the cosine similarity between the centroids of task $i$ and task $j$.
    \item Perform the Spectral (Eigenvalue) Decomposition of the overlap matrix $S$:
    \begin{equation}
        S = V \Lambda V^T
    \end{equation}
    where $V \in \mathbb{R}^{K \times K}$ is an orthonormal matrix of eigenvectors, and $\Lambda = \text{diag}(\lambda_1, \lambda_2, \dots, \lambda_K)$ is a diagonal matrix containing the positive eigenvalues.
    \item Compute the symmetric inverse square root matrix $S^{-1/2}$:
    \begin{equation}
        S^{-1/2} = V \Lambda^{-1/2} V^T = V \text{diag}\left(\lambda_1^{-1/2}, \lambda_2^{-1/2}, \dots, \lambda_K^{-1/2}\right) V^T
    \end{equation}
    \item Project the original centroids using the symmetric de-entangling operator:
    \begin{equation}
        C_{\text{SMD}} = S^{-1/2} C
    \end{equation}
\end{enumerate}

We mathematically verify that $C_{\text{SMD}}$ is orthonormal:
\begin{equation}
    C_{\text{SMD}} C_{\text{SMD}}^T = (S^{-1/2} C) (S^{-1/2} C)^T = S^{-1/2} (C C^T) S^{-1/2} = S^{-1/2} S S^{-1/2} = I_{K \times K}
\end{equation}
Because the transformation matrix $S^{-1/2}$ is symmetric, SMD applies an identical, unbiased orthogonalization force to all centroids. This preserves the symmetric structure of the representational manifold while cleanly de-entangling the task templates, suppressing routing flicker under high representation noise.

\subsection{Online Expectation-Maximization (EM) Updates for GMM Adaptation}
\label{sec:appendix_em}

The diagonal Coordinate Gaussian Mixture Model (GMM) safety shield fits a probability density over the similarity coordinates $u'_b \in \mathbb{R}^K$. To handle long-term environmental, illumination, and covariate shifts on physical edge devices without triggering costly offline recalibration, we formulate an on-device online EM parameter adaptation loop.

Let $\theta^{(t)} = \{\pi_c^{(t)}, \mu_c^{(t)}, \Sigma_c^{(t)}\}_{c=1}^C$ denote the GMM parameters at time step $t$, where $\pi_c$ is the mixing coefficient, $\mu_c \in \mathbb{R}^K$ is the mean, and $\Sigma_c = \text{diag}(\sigma_{c, 1}^2, \dots, \sigma_{c, K}^2)$ is the diagonal covariance matrix of component $c$. For each incoming similarity coordinate vector $u'_b$ that is accepted as in-distribution by the GMM shield, we dynamically update the parameters online:

\begin{enumerate}
    \item \textbf{Expectation Step (E-step):} Compute the posterior probability (responsibility) $\gamma_{c}(u'_b)$ of component $c$ generating the coordinate $u'_b$:
    \begin{equation}
        \gamma_{c}(u'_b) = \frac{\pi_c^{(t)} \mathcal{N}\left(u'_b \;\middle|\; \mu_c^{(t)}, \Sigma_c^{(t)}\right)}{\sum_{j=1}^C \pi_j^{(t)} \mathcal{N}\left(u'_b \;\middle|\; \mu_j^{(t)}, \Sigma_j^{(t)}\right)}
    \end{equation}
    \item \textbf{Maximization Step (M-step):} Update the parameters using an exponentially weighted moving average (EWMA) with a slow learning rate $\alpha_{\text{EM}} \in [0.01, 0.05]$:
    \begin{align}
        \pi_c^{(t+1)} &= (1 - \alpha_{\text{EM}}) \pi_c^{(t)} + \alpha_{\text{EM}} \gamma_{c}(u'_b) \\
        \mu_c^{(t+1)} &= (1 - \alpha_{\text{EM}} \eta_c) \mu_c^{(t)} + \alpha_{\text{EM}} \eta_c u'_b \\
        \sigma_{c, k}^{2 (t+1)} &= (1 - \alpha_{\text{EM}} \eta_c) \sigma_{c, k}^{2 (t)} + \alpha_{\text{EM}} \eta_c \left(u'_{b, k} - \mu_{c, k}^{(t+1)}\right)^2
    \end{align}
    where $\eta_c = \gamma_{c}(u'_b) / \pi_c^{(t+1)}$ represents the normalized learning rate scaling factor for component $c$.
\end{enumerate}
Because our similarity coordinate space is extremely low-dimensional ($K=4$), executing this online GMM update requires less than 100 Floating-Point Operations (FLOPs) per sample, which is several orders of magnitude smaller than a single backbone block execution. This allows GMM updates to run seamlessly in a background thread of a mobile application without causing any pipeline latency or thread-synchronization stalls.

\section{Experimental Specifications and Hyperparameters}
\label{sec:appendix_params}

This section details the hyperparameter configurations, dataset statistics, and hardware-modeling parameters utilized across our analytical simulations and physical weight-quantization validations.

\subsection{Isolating Coordinate Sandbox (ICS) Setup}
\label{sec:appendix_ics_params}

The ICS sandbox simulates a 12-layer Vision Transformer backbone (ViT-Tiny) with $K=4$ registered specialized experts corresponding to MNIST, Fashion-MNIST, CIFAR-10, and SVHN.
\begin{itemize}
    \item \textbf{Backbone Dimension ($D$):} 192 (hidden state dimension of ViT-Tiny).
    \item \textbf{Expert Adapter Rank ($r$):} 8 (rank of LoRA matrices $A_k, B_k$).
    \item \textbf{Subspace Dimension ($D_{\text{coord}}$):} 48-dimensional orthogonal blocks of Layer 3 representations are mapped to task similarity coordinates, establishing the idealized boundary condition used to isolate systems variables.
    \item \textbf{Routing Layer Index:} Layer 3 (ZCA-IDC is executed on the outputs of the third transformer block).
    \item \textbf{Routing Temperature ($\tau$):} 0.001 (extremely sharp, approximating one-hot expert routing).
    \item \textbf{Expert Gating Threshold ($\theta$):} 0.01 (expert pathways with routing coefficients $w_{k, b} < 0.01$ are bypassed losslessly).
    \item \textbf{Representation Noise ($\sigma_{\text{noise}}$):} Varying in $[0.0, 0.6]$ (standard Gaussian noise added to Layer 3 representations to evaluate router stability).
    \item \textbf{Centroid Entanglement Factor ($\epsilon$):} Varying in $[0.0, 0.9]$ (controls the overlap and cosine similarity between task centroids to evaluate de-entangling orthogonalizations).
\end{itemize}

\subsection{Real pre-trained Weights Quantization Validation}
\label{sec:appendix_pretrained_params}

For our empirical post-training quantization calibration checks on real Vision Transformer weights, we adopt the following configuration:
\begin{itemize}
    \item \textbf{Pre-trained Model:} \texttt{vit\_tiny\_patch16\_224} (loaded from the \texttt{timm} library, containing 5.7M parameters pre-trained on ImageNet-1k).
    \item \textbf{Target Layers for Validation:} The MLP \texttt{fc1} layers of Blocks 5, 9, and 12 (capturing early, middle, and late representation depth profiles of the ViT backbone).
    \item \textbf{LoRA Construction:} Constructed via Singular Value Decomposition (SVD) on the pre-trained linear projection weights of the respective blocks, extracting the top-$r$ singular vectors to form low-rank adapters with rank $r=8$.
    \item \textbf{Evaluation Dataset:} The CIFAR-10 test set (loaded via \texttt{torchvision}). We propagate a batch of 256 real images through the early blocks of the pre-trained ViT to extract 50,176 real visual activation token vectors at Layer 4.
    \item \textbf{Calibration Split:} 10,000 activation tokens (used as the representative calibration subset to optimize QASC scale factors offline).
    \item \textbf{Test Split:} 40,176 activation tokens (held-out subset used to report reconstruction Mean Squared Error and Cosine Similarity).
\end{itemize}

\subsection{Hardware Modeling and Systems Cost Parameters}
\label{sec:appendix_hardware_params}

Our hardware-aware execution latency model represents the CPU and cache hierarchy of a standard quad-core ARM Cortex-A72 processor on a Raspberry Pi 4 board.
\begin{itemize}
    \item \textbf{Base Model Size ($M_{\text{base}}$):} 22.8 MB (5.7M parameters stored in FP32 format).
    \item \textbf{FP32 Expert Size ($M_{\text{LoRA\_FP32}}$):} 2.76 MB (combined weight size of 4 task-specific experts with rank $r=8$).
    \item \textbf{INT8 Expert Size ($M_{\text{LoRA\_INT8}}$):} 0.69 MB (combined expert weight size under 8-bit quantization).
    \item \textbf{INT4 Expert Size ($M_{\text{LoRA\_INT4}}$):} 0.345 MB (combined expert weight size under 4-bit uniform quantization).
    \item \textbf{Peak DRAM Bandwidth ($B_{\text{mem}}$):} 4.4 GB/s (representing LPDDR4 memory bandwidth constraints on Pi 4).
    \item \textbf{L1/L2 Cache Capacity:} 32 KB L1 Instruction/Data caches per core, and a shared 1 MB L2 cache.
    \item \textbf{Thread Synchronization Barrier ($T_{\text{sync}}$):} 0.5 ms (representing the context-switching and mutex scheduling overhead to coordinate multi-threaded jobs across CPU cores).
    \item \textbf{Inter-Cluster Cross Penalty ($T_{\text{cross-cluster}}$):} 0.15 ms (representing the inter-cluster cache coherence penalty to offload task queue jobs from Big cores to LITTLE cores).
    \item \textbf{Unpacking/Casting Compute Penalty:} 15\% latency overhead added to INT4 matrix multiplication operations to model the instruction-level register unpacking (bit-shifting and masking) penalty.
\end{itemize}

\section{Additional Systems and Empirical Profiling}
\label{sec:appendix_additional_results}

This section provides additional quantitative evaluations and systems profiling to supplement the experimental results presented in Section~4.

\subsection{Quantization Scale Optimization and Discretization Error Profiles}
\label{sec:appendix_discretization_profiles}

To illustrate how our proposed post-training Quantization-Aware Scale Calibration (QASC) minimizes discretization rounding error and activation clipping noise, we profile the Mean Squared Error (MSE) surface of the low-rank up-projection weight matrix $B_{12} \in \mathbb{R}^{192 \times 8}$ under 4-bit uniform quantization.

Table~\ref{tab:qasc_grid_mse} compares the reconstruction error and output representation fidelity under three different quantization configurations: (1) standard uncalibrated Round-To-Nearest (RTN), (2) traditional MinMax dynamic range clipping, and (3) our sequentially decoupled QASC scale calibration.

\begin{table}[h]
    \caption{Quantization Reconstruction Error and Representation Fidelity Profiles for Up-Projection Weight Block $B_{12}$ under 4-Bit Uniform Quantization evaluated over 40,176 CIFAR-10 Visual Tokens.}
    \label{tab:qasc_grid_mse}
    \vskip 0.15in
    \begin{center}
    \begin{small}
    \begin{sc}
    \begin{tabular}{lccc}
        \toprule
        Quantization Scheme & Scale Factor $s_B$ & Reconstruction MSE & Output Cosine Similarity \\
        \midrule
        Unquantized (FP32) & -- & 0.0000 & 1.0000 \\
        RTN (Standard PTQ) & 0.0524 & 0.0824 & 0.9416 \\
        MinMax Clipping     & 0.0431 & 0.0512 & 0.9682 \\
        \textbf{QASC (Ours)} & \textbf{0.0384} & \textbf{0.0214} & \textbf{0.9894} \\
        \bottomrule
    \end{tabular}
    \end{sc}
    \end{small}
    \end{center}
    \vskip -0.1in
\end{table}

As shown in Table~\ref{tab:qasc_grid_mse}, standard Round-To-Nearest (RTN) suffers from extreme discretization rounding error because it sets the scale factor based on the absolute maximum value of the weight tensor ($\max(|B_{12}|)$), which is heavily biased by outliers. Traditional MinMax clipping improves upon RTN by slightly shrinking the dynamic range boundary (clipping extreme outliers to reduce rounding step size for the majority of entries), but its heuristic search is sub-optimal.

In contrast, our proposed QASC scale calibration sequentially optimizes the scale factor to minimize the reconstruction error under the pre-quantized input activation path. By finding the optimal scale factor $s_B^* = 0.0384$ (which slightly clips the top 1.5\% of weight outliers to drastically reduce the uniform rounding step size for the remaining 98.5\% of weights), QASC slashes the reconstruction Mean Squared Error by over \textbf{74\%} compared to RTN, recovering an outstanding output representation Cosine Similarity of \textbf{0.9894}.

\subsection{Energy and Power Consumption Optimization Profiles}
\label{sec:appendix_energy_profiles}

To ground our systems co-design claims, we model and report the energy consumption profiles of our framework on physical edge CPU constraints. Under physical serving, the energy consumed per inference batch $E_{\text{batch}}$ consists of: (1) the active processor compute energy $E_{\text{compute}}$, and (2) the DRAM-to-SRAM weight transfer energy $E_{\text{transfer}}$:
\begin{equation}
    E_{\text{batch}} = E_{\text{compute}} + E_{\text{transfer}} = \left( P_{\text{CPU\_active}} \times T_{\text{compute}} \right) + \left( E_{\text{DRAM\_access}} \times M_{\text{DRAM\_transfer}} \right)
\end{equation}
We calibrate this model using empirical physical statistics from the Raspberry Pi 4 CPU and DRAM: peak active CPU core power $P_{\text{CPU\_active}} \approx 2.1$ Watts, and DRAM access energy $E_{\text{DRAM\_access}} \approx 65.0$ pJ/bit (or $0.52$ milli-Joules per Megabyte of transferred data).

Table~\ref{tab:energy_consumption_profile} reports the dynamic energy consumption per batch ($B=256$) and cumulative energy over 1,024 samples under highly task-interleaved heterogeneous streaming demands.

\begin{table}[h]
    \caption{Projected Edge Hardware Energy Consumption Profile under Highly Task-Interleaved Heterogeneous Streams (1024 Samples total).}
    \label{tab:energy_consumption_profile}
    \vskip 0.15in
    \begin{center}
    \begin{small}
    \begin{sc}
    \begin{tabular}{lcccc}
        \toprule
        Method & Weight Precision & Compute Energy (J) & Transfer Energy (mJ) & Cumulative Energy (J) \\
        \midrule
        PFSR + MBH SOTA & FP32 & 1.5746 J & 48.86 mJ & 1.6234 J \\
        SPS-ZCA         & FP32 & 0.4166 J & 13.29 mJ & 0.4299 J \\
        Q-SPS (Ours)    & INT8 & 0.4059 J & 12.21 mJ & 0.4181 J \\
        \textbf{CG-Q-SPS (Ours)} & \textbf{INT4} & \textbf{0.3412 J} & \textbf{6.14 mJ} & \textbf{0.3473 J} \\
        \bottomrule
    \end{tabular}
    \end{sc}
    \end{small}
    \end{center}
    \vskip -0.1in
\end{table}

As shown in Table~\ref{tab:energy_consumption_profile}, the standard SOTA PFSR+MBH consumes an elevated \textbf{1.623 J} of cumulative energy because it must run $G=4$ sequential passes of the base model over split sub-batches to prevent activation bleeding, multiplying active processor compute energy.

By ensembling sample-wise in a single parallel pass, our unquantized baseline SPS-ZCA slashes cumulative serving energy by over \textbf{73.5\%} (down to 0.429 J). When we apply low-bit uniform weight quantization and our conditional execution gating, \textbf{CG-Q-SPS (INT4)} achieves the ultimate energy frontier, consuming only \textbf{0.347 J} of cumulative energy. This represents an exceptional \textbf{78.6\% energy reduction} compared to PFSR+MBH SOTA and a \textbf{19.2\% reduction} compared to SPS-ZCA, demonstrating that compressing LoRA weights to INT4 and bypassing inactive experts dynamically is a highly effective co-design strategy to extend the battery lifespan of mobile and embedded IoT nodes.
